For the paired region analysis, we searched all 2{,}838 qualified dark models and 3{,}013
qualified controls as complete chains and as confidence-defined regions against the frozen
PDB Foldseek database downloaded on 30 July 2026. A centred 15-residue moving average of
per-residue pLDDT, shortened at chain ends, identified intervals at or above 0.70.
We bridged gaps of at most five residues when their mean unsmoothed pLDDT was at least 0.55,
then retained intervals of at least 50 residues. Each model contributed a positive crop outcome
when any eligible region reached query-normalised TM-score 0.5. Models without eligible
regions remained in the unconditional denominator. Both search arms used Foldseek
\texttt{10-941cd33}, \texttt{--alignment-type 1 -e 10 -s 9.5 --max-seqs 1000}.
Queries reaching the hit cap while remaining unmatched under any exclusion tier were rerun
with \texttt{--max-seqs 5000}. The primary tier removed each control's source PDB entry;
further tiers removed alignments at or above 30\%, 25\% and 20\% sequence identity.
The GPU sensitivity used the same structures and padded frozen reference on ORIX, with
\texttt{--gpu 1 --prefilter-mode 1 --max-seqs 1000}. This changes the retrieval algorithm to
ungapped prefiltering. Supplementary Note S3 reports cap saturation and CPU--GPU disagreement.

The paired analysis uses query-normalised scores throughout. An audit of the Foldseek build's
alignment-normalised score found that its normalisation length omitted one residue from
inclusive alignment spans. The same length also enters the distance scale and optimisation;
exact correction requires rescoring the coordinates. Raw alignment-normalised values are
preserved as software diagnostics. The whole-chain and crop endpoints and the pooled
1{,}031-model dark group depend exclusively on query-normalised scores, whose normalisation
uses the full query length in this implementation.
